% !TeX root = ../main.tex
\section{Impact of Research}

This research has the potential to deliver impact by making soccer analytics more accessible, more energy-efficient, and less dependent on expensive cloud-based or specialist commercial infrastructure. While the work presented in this thesis is still at a prototype and benchmarking stage, it demonstrates that useful computer-vision-driven match analysis can be performed on an embedded edge device, which creates a realistic path toward lower-cost deployment in coaching, analysis, and broadcasting settings.

\subsection{Potential Beneficiaries and Pathways to Impact}

The most immediate beneficiaries of this work are researchers and engineers working on embedded computer vision, since the thesis provides a structured comparison of deployment-focused optimisation methods across object detection, pitch keypoint localisation, and full-stack soccer analysis. The results also have potential value for sports science groups, university teams, amateur and semi-professional clubs, and smaller broadcasters, all of whom may lack the budget required for commercial analytics providers or high-end cloud infrastructure.

In the short term, the main impact of the project is methodological. The benchmarking framework, the compression analysis, and the integrated pipeline provide evidence for which optimisation strategies are most effective when the goal is to balance accuracy with real-time or near-real-time operation on constrained hardware. This can inform future academic work and help guide deployment decisions for similar sports analytics systems.

In the medium term, this research could support the development of low-cost tools for match review, player development, and local broadcasting. For example, a refined version of the system could help coaches and analysts generate statistics such as possession trends, pass patterns, and player movement indicators without the need for large-scale cloud processing. This could widen access to data-driven decision-making beyond elite clubs and make analytical tools available to organisations with more limited resources.

In the longer term, if the methods were extended into a more robust production system, the work could contribute economically by supporting innovation in sports technology products and services. This includes possible applications in affordable analytics platforms, edge-deployed performance tools, and broadcast support systems. The impact would not arise from the prototype alone, but from the adoption and refinement of its outputs into deployable products, workflows, or operational standards.

\subsection{Ethical Considerations and Risk Mitigation}

Although the project is motivated by performance analysis rather than surveillance, it still raises important ethical considerations. The first concerns the use of player image data in training and evaluation datasets. In this work, the datasets were drawn from publicly available or appropriately shared football video sources used for non-commercial research purposes. In addition, the system does not perform facial recognition, biometric identification, or any attempt to infer personal identity. Players are represented only through bounding boxes, keypoints, and team-level analytical outputs, which reduces the privacy sensitivity of the system compared with person-identification technologies.

However, some risks remain. A system of this type could be used inappropriately for unauthorised monitoring of athletes, or its outputs could be over-interpreted in settings where model predictions are uncertain. Dataset bias is also a concern, since models trained primarily on professional broadcast footage may not generalise equally well to lower-quality video, different lighting conditions, youth matches, or non-broadcast camera angles. If such a system were used without caution, inaccurate detections or biased model behaviour could influence tactical decisions, player evaluation, or analyst confidence in misleading ways.

These risks can be mitigated in several ways. First, the intended use of the system should remain limited to performance analysis rather than identity-based tracking. Second, model outputs should be treated as decision-support tools rather than fully autonomous judgements, with human review maintained for important tactical or evaluative decisions. Third, deployment should include transparent reporting of uncertainty, known failure cases, and data limitations. Finally, any future expansion toward commercial or operational use should ensure that dataset licensing, consent requirements, and governance procedures are clearly respected.

\subsection{Sustainability Considerations and Sustainable Development Goals}

This project is also relevant to sustainability because it focuses on computational efficiency rather than maximising performance through ever-larger models and unrestricted compute usage. Deep learning systems can carry a significant environmental cost when they rely on repeated training cycles or cloud-scale inference without regard for efficiency \cite{desislavov2023trends}. By evaluating hardware acceleration, quantisation, pruning, and knowledge distillation for embedded deployment, this thesis instead promotes an approach in which useful analytical performance is achieved with lower power draw, reduced thermal load, and less dependence on continuous data transfer to centralised servers.

The work aligns most directly with United Nations Sustainable Development Goal 9, \textit{Industry, Innovation and Infrastructure}, by contributing toward practical and affordable digital infrastructure for applied sports analytics. It also aligns with Sustainable Development Goal 12, \textit{Responsible Consumption and Production}, by emphasising efficient model design and resource-aware deployment rather than unnecessary compute expansion. In addition, it has relevance to Sustainable Development Goal 13, \textit{Climate Action}, since edge-based inference and model compression can reduce the energy overhead associated with cloud processing and prolonged high-power operation.

At the same time, sustainability is not automatically guaranteed simply because inference is performed on-device. Training many model variants can still consume substantial energy, and future development could increase hardware turnover or encourage inefficient experimentation if left unmanaged. For that reason, this project has aimed to limit redundant training runs, reuse model outputs where possible, and prioritise controlled benchmarking. A responsible continuation of this research would continue to favour energy-aware evaluation, selective experimentation, and deployment choices that are justified by clear analytical benefit rather than by marginal performance gains alone.
